\documentclass[10pt]{article}
\usepackage[T1]{fontenc}
\usepackage[utf8]{inputenc}
\usepackage[french]{babel}
\usepackage{amsmath,amssymb,mathrsfs,mathpazo}
\usepackage{geometry}
\geometry{a4paper,margin=2cm}
\usepackage{graphicx,array,tabularx,tabularray,colortbl,multirow,multicol,booktabs,enumitem,pifont,ulem,fancybox,diagbox,cancel}
\usepackage{pstricks,pstricks-add,pst-plot,pst-tree,pst-node,pst-text,pst-func,pst-math,pst-bspline,pst-3dplot,pst-eucl,pst-solides3d}
\let\pscalculate\undefined
\usepackage{pst-bezier}
\usepackage{tikz,pgfplots,tkz-tab}
\usetikzlibrary{arrows,arrows.meta,patterns,calc,positioning,trees,decorations.markings,shapes.geometric,shapes.misc,shapes.arrows,intersections,decorations.pathmorphing,decorations.pathreplacing,backgrounds,shadows,babel}
\pgfplotsset{compat=1.18}
\usepgfplotslibrary{fillbetween}
\usepackage[np]{numprint}
\usepackage{listings,makecell,pas-tableur,twoopt}
\let\barre\undefined
\usepackage[pstricks]{tablvar}
\definecolor{darkgreen}{RGB}{0,100,0}
\providecommand{\red}{\color{red}}
\providecommand{\blue}{\color{blue}}
\providecommand{\green}{\color{green}}
\providecommand{\black}{\color{black}}
\providecommand{\white}{\color{white}}
\providecommand{\gray}{\color{gray}}
\providecommand{\cyan}{\color{cyan}}
\providecommand{\magenta}{\color{magenta}}
\providecommand{\yellow}{\color{yellow}}
\providecommand{\violet}{\color{violet}}
\providecommand{\orange}{\color{orange}}
\providecommand{\euro}{\texteuro}
\providecommand{\EUR}{\texteuro}
\providecommand{\ieme}{\textsuperscript{e}}
\providecommand{\degres}{\ensuremath{^\circ}}
\providecommand{\degree}{\ensuremath{^\circ}}
\providecommand{\vv}[1]{\overrightarrow{#1}}
\providecommand{\Coord}[1]{#1}
\providecommand{\decosix}{}
\providecommand{\starredbullet}{\ensuremath{\bullet}}
\pagestyle{empty}
\setlength{\parindent}{0pt}
\usepackage{adjustbox}\DeclareUnicodeCharacter{2236}{:}\DeclareUnicodeCharacter{2190}{\ensuremath{\leftarrow}}
\providecommand{\R}{\mathbb{R}}
\providecommand{\N}{\mathbb{N}}
\providecommand{\Z}{\mathbb{Z}}
\providecommand{\Q}{\mathbb{Q}}
\providecommand{\D}{\mathbb{D}}
\providecommand{\C}{\mathbb{C}}
\providecommand{\e}{{\mathrm{e}}}
\providecommand{\vect}[1]{\overrightarrow{#1}}
\providecommand{\vectt}[1]{\overrightarrow{\mathrm{#1}}}
\providecommand{\vvt}[1]{\overrightarrow{\mathrm{#1}}}
\providecommand{\np}[2][]{{#2~#1}}
\providecommand{\vv}[1]{{\overrightarrow{#1}}}
\renewcommand{\d}{{\mathrm{d}}}
\providecommand{\degres}{{^\circ}}
\providecommand{\degre}{{^\circ}}
\providecommand{\blue}{\color{blue}}
\providecommand{\red}{\color{red}}
\providecommand{\enskip}{\quad}
\providecommand{\llbracket}{\lbrack\!\lbrack}
\providecommand{\rrbracket}{\rbrack\!\rbrack}
\providecommand{\euro}{{\text{€}}}
\providecommand{\no}{{\text{n°}}}
\providecommand{\interval}[2]{{[#1;#2]}}
\providecommand{\intervFO}[2]{{[#1;#2[}}
\providecommand{\dotfill}{\ldots}

\providecommand{\ds}{\displaystyle}
\providecommand{\textding}[1]{#1}
\providecommand{\Oij}{\(\mathrm{O};\vec{\imath},\vec{\jmath}\)}
\providecommand{\Oijk}{\(\mathrm{O};\vec{\imath},\vec{\jmath},\vec{k}\)}
\providecommand{\Ouv}{\(\mathrm{O};\vec{u},\vec{v}\)}
\providecommand{\intervFF}[2]{\ensuremath{[#1;#2]}}
\renewcommand{\euro}{\ifmmode\text{\texteuro}\else\texteuro\fi}
\DeclareUnicodeCharacter{00B0}{\ifmmode^{\circ}\else\textdegree\fi}
\DeclareUnicodeCharacter{20AC}{\ifmmode\text{\texteuro}\else\texteuro\fi}
% Énoncés : Éducation nationale. Transcription et corrigés : APMEP, Denis Vergès et contributeurs. Découpage et classement : CoopMaths, Sébastien Lozano, Éric Elter et contributeurs. Adaptation de la présentation : Maths Explorer.
% CC BY-SA : https://coopmaths.fr/www/about/
% https://coopmaths.fr/alea/static/e3c/2020/tex/e3c_2020_00_sujet10_2.tex
\begin{document}
\noindent\textit{APMEP / CoopMaths. CC BY-SA. Présentation adaptée par Maths Explorer.}\par\medskip
\medskip

Un industriel souhaite fabriquer une boîte sans couvercle à partir d'une plaque de métal de
$18$~cm de largeur et de $24$~cm de longueur. Pour cela, il enlève des carrés dont la longueur du côté mesure $x$ cm aux quatre coins de la pièce de métal et relève ensuite verticalement pour fermer les côtés.

\begin{center}
\psset{unit=0.8cm}
\begin{pspicture}(16,7)

\psframe(1,1)(7,5)
\pspolygon[linestyle=dashed,fillstyle=solid,fillcolor=lightgray](1,2)(1,4)(2,4)(2,5)(6,5)(6,4)(7,4)(7,2)(6,2)(6,1)(2,1)(2,2)
\psline{<->}(6,0.5)(7,0.5)\uput[d](6.5,0.5){$x$}
\psline{<->}(7.4,1)(7.4,2)\uput[r](7.4,1.5){$x$}
\psline{<->}(0.3,1)(0.3,5)\psline{<->}(1,5.6)(7,5.6)
\uput[l](0.3,3){18}\uput[u](4,5.6){24}
\psframe[fillstyle=solid,fillcolor=lightgray](10,1)(14,2)
\pspolygon[fillstyle=solid,fillcolor=lightgray](14,1)(16,2.8)(16,3.8)(14,2)
\psline(10,2)(12,3.8)(16,3.8)
\psline(12,3.8)(12,2.8)(14.85,2.8)
\psline(11,2)(12,2.8)
\psarc[linewidth=0.15cm]{<-}(8.5,0){3}{70}{110}
\end{pspicture}
\end{center}

Le volume de la boîte ainsi obtenue est une fonction définie sur l'intervalle [0~;~9] notée
$V(x)$.

\medskip

\begin{enumerate}
\item Justifier que pour tout réel $x$ appartenant à [0~;~9] : $V(x) = 4x^3 - 84x^2 + 432x$. 
\item On note $V'$ la fonction dérivée de $V$ sur [0~;~9]. \\
Donner l'expression de $V'(x)$ en fonction de $x$.
\item Dresser alors le tableau de variations de $V$ en détaillant la démarche.
\item Pour quelle(s) valeur(s) de $x$ la contenance de la boîte est-elle maximale ?
\item L'industriel peut-il construire ainsi une boîte dont la contenance est supérieure ou égale à $650$ cm$^3$ ? Justifier.
\end{enumerate}
\bigskip

\end{document}
